% ============================================================
%  Sistema Métrico — Apostila Premium | PratiqueJá
%  Engine: XeLaTeX
% ============================================================
\documentclass[12pt]{article}
\usepackage{../../pratiqueja}
\usepackage{tikz}
\usetikzlibrary{arrows.meta}

% \hlm — destaque de resultado em modo-math (cor sem bold)
\newcommand{\hlm}[1]{{\color{darkblue}#1}}

\tikzset{
  deg/.style={draw=accentblue!60!black, line width=0.8pt, fill=accentblue!12!white, rounded corners=2pt},
  degb/.style={draw=darkblue, line width=1pt, fill=rulebg, rounded corners=2pt},
}

\setsubject{Sistema Métrico}

\begin{document}
\pagenumbering{arabic}
\setstretch{1.2}
\color{bodytext}

\heroheader{Sistema Métrico}%
           {Conversão de Unidades de Medida}%
           {Matemática · Básico}

\setlength{\columnsep}{18pt}
\setlength{\columnseprule}{0.5pt}
\def\columnseprulecolor{\color{exborder}}

\begin{multicols}{2}

%─────────────────────────────────────────────────────────────
\TW{m}{badgepurple}{Comprimento}{Unidade base: metro (m)}

\regra{A unidade padrão é o \textbf{metro (m)}. As demais são
múltiplos ou submúltiplos, separados por fatores de $10$.}

\SH{A escada das unidades}
\begin{center}\vspace{2pt}
\begin{tikzpicture}[scale=0.95]
\foreach \i/\u in {0/km,1/hm,2/dam} {
  \fill[deg] (\i*1.0,-\i*0.52) rectangle ++(0.95,0.46);
  \node[font=\footnotesize\bfseries, accentblue!50!black]
        at (\i*1.0+0.47,-\i*0.52+0.23) {\u};
}
\fill[degb] (3*1.0,-3*0.52) rectangle ++(0.95,0.46);
\node[font=\footnotesize\bfseries, darkblue] at (3*1.0+0.47,-3*0.52+0.23) {m};
\foreach \i/\u in {4/dm,5/cm,6/mm} {
  \fill[deg] (\i*1.0,-\i*0.52) rectangle ++(0.95,0.46);
  \node[font=\footnotesize\bfseries, accentblue!50!black]
        at (\i*1.0+0.47,-\i*0.52+0.23) {\u};
}
\node[font=\scriptsize, vegreen!50!black] at (1.4,0.55) {$\times 10$ ←};
\node[font=\scriptsize, vered!80!black] at (5.0,-3.35) {→ $\div 10$};
\end{tikzpicture}
\end{center}

\obs{Cada passo à \textbf{direita}: $\div10$. \quad Cada passo à
\textbf{esquerda}: $\times10$.}

\exemp{%
  $3{,}5\,\text{km} = \hlm{3500\,\text{m}}$ \;{\footnotesize($\times1000$)}\\[3pt]
  $250\,\text{cm} = \hlm{2{,}5\,\text{m}}$ \;{\footnotesize($\div100$)}\\[3pt]
  $45\,\text{mm} = \hlm{4{,}5\,\text{cm}}$ \;{\footnotesize($\div10$)}%
}

%─────────────────────────────────────────────────────────────
\TW{kg}{darkblue}{Massa}{Unidade base: quilograma (kg)}

\regra{A unidade do SI é o \textbf{quilograma (kg)}, mas o
\emph{grama (g)} é o centro das conversões na escala métrica.}

\formula{t \;→\; kg \;→\; hg \;→\; dag \;→\; g \;→\; dg \;→\; cg \;→\; mg}

\exemp{%
  $1\,\text{t} = 1000\,\text{kg}$ \qquad $1\,\text{kg} = 1000\,\text{g}$\\[3pt]
  $2{,}5\,\text{kg} = \hlm{2500\,\text{g}}$\\[3pt]
  $350\,\text{g} = \hlm{0{,}35\,\text{kg}}$\\[3pt]
  $1{,}5\,\text{t} = \hlm{1500\,\text{kg}}$%
}

%─────────────────────────────────────────────────────────────
\TW{L}{darkblue}{Capacidade}{Unidade base: litro (L)}

\regra{A unidade é o \textbf{litro (L)}; as conversões seguem o mesmo
padrão de múltiplos e submúltiplos de $10$.}

\formula{kL \;→\; hL \;→\; daL \;→\; L \;→\; dL \;→\; cL \;→\; mL}

\exemp{%
  $1\,\text{L} = 1000\,\text{mL}$ \qquad $1\,\text{kL} = 1000\,\text{L}$\\[3pt]
  $2\,\text{L} = \hlm{2000\,\text{mL}}$\\[3pt]
  $500\,\text{mL} = \hlm{0{,}5\,\text{L}}$\\[3pt]
  $3{,}5\,\text{kL} = \hlm{3500\,\text{L}}$%
}

%─────────────────────────────────────────────────────────────
\TW{m²/m³}{badgegreen}{Área e Volume}{Conversões quadráticas e cúbicas}

\regra{\textbf{Área:} cada passo na escala de comprimento vale
$\times100$ ($10^2$).\\[2pt]
\textbf{Volume:} cada passo vale $\times1000$ ($10^3$).}

\SH{Área (passo $\times100$ / $\div100$)}
\exemp{%
  km² → hm² → dam² → m² → dm² → cm² → mm²\\[2pt]
  $1\,\text{m}^2 = \hlm{10\,000\,\text{cm}^2}$\\[3pt]
  $1\,\text{km}^2 = \hlm{1\,000\,000\,\text{m}^2}$%
}

\SH{Volume (passo $\times1000$ / $\div1000$)}
\exemp{%
  $1\,\text{m}^3 = \hlm{1\,000\,000\,\text{cm}^3}$\\[3pt]
  $1\,\text{dm}^3 = \hlm{1\,\text{L}}$ \;{\footnotesize(elo com a capacidade)}%
}

%─────────────────────────────────────────────────────────────
\TW{h}{badgegreen}{Tempo}{Horas, minutos e segundos}

\nota{O tempo \textbf{não} é decimal: usa fatores fixos de $60$
(entre s, min e h) e $24$ (horas no dia).}

\formula{$1\,\text{h} = 60\,\text{min}$ \quad $1\,\text{min} = 60\,\text{s}$\\[3pt]
$1\,\text{h} = 3600\,\text{s}$ \quad $1\,\text{dia} = 24\,\text{h}$}

\exemp{%
  $2{,}5\,\text{h} = 2\,\text{h}\,30\,\text{min} = \hlm{150\,\text{min}}$\\[3pt]
  $90\,\text{min} = 1{,}5\,\text{h} = \hlm{1\,\text{h}\,30\,\text{min}}$\\[3pt]
  $7200\,\text{s} = 120\,\text{min} = \hlm{2\,\text{h}}$%
}

\nota{\textbf{Cuidado:} $2{,}5\,\text{h}$ \emph{não} é
$2\,\text{h}\,50\,\text{min}$. A parte decimal é fração de hora:
$0{,}5\,\text{h} = 30\,\text{min}$.}

\end{multicols}

\end{document}
